% !TEX root =  Lecture6.tex


%%%%%%%%%%%%%%%%%%%%%%%%%%%%%%%%%%%%
% Recap/Agenda
%%%%%%%%%%%%%%%%%%%%%%%%%%%%%%%%%%%%
\begin{frame}
\frametitle{Module 1: Modeling and Linear Regression}
    \begin{itemize}
    \item \hl{Previously:} multiple regression --- conditional interpretation, indicator variables, and adjusted $R^2$.
    \item \hl{This time:} logistic regression, when the response is a \emph{category} rather than a number.
      \begin{itemize}
      \item the logistic regression model for a binary response
      \item odds, log odds, and coefficients as odds ratios
      \item fitting by maximum likelihood
      \item classification thresholds and unbalanced data
      \end{itemize}
    \item \hl{Reading:} IMS Chapter 9.1--9.2 and 9.4.
    \item \hl{Homework for today:}
    \item \hl{Homework for next class:}
    \end{itemize}
\end{frame}


%%%%%%%%%%%%%%%%%%%%%%%%%%%%%%%%%%%%
% Learning objectives:
%%%%%%%%%%%%%%%%%%%%%%%%%%%%%%%%%%%%
\begin{frame}
    \frametitle{Lecture Objectives}
    \goals{%
    \begin{itemize}
    \item explain the logistic regression model and justify its use for a binary outcome;
    \item convert between probability, odds, and log odds;
    \item interpret a coefficient as an \emph{odds ratio}, $e^{\beta_j}$;
    \item describe how the model is fit by maximum likelihood;
    \item choose a classification threshold, and explain why accuracy is misleading on unbalanced data.
    \end{itemize}}
    \objectives{%
    \begin{itemize}
    \item \textbf{(M1, LO6)} Logistic Regression Models
    \end{itemize}}
\end{frame}
